% --- Archon physics formalization source begin ---
% archon:physics
% archon:covers PhyXMiniProblems/problem_phyx_mini_0915.lean
% archon:source-report reports/phyx_mini/problem_phyx_mini_0915.source.json

\chapter{Physics problem phyx\_mini\_0915}
\label{ch:PhyXMiniProblems_problem_phyx_mini_0915}

\paragraph{Problem source.}
\#\# Question

What is the strength of the electric field at the position indicated by the dot?

\#\# Answer choices

- A: A: 1.3 \textbackslash{}times 10\textasciicircum{}\{3\}N/C

- B: B: 5.3 \textbackslash{}times 10\textasciicircum{}\{3\}N/C

- C: C: 1.0 \textbackslash{}times 10\textasciicircum{}\{4\}N/C

- D: D: 3.5 \textbackslash{}times 10\textasciicircum{}\{3\}N/C

\#\# Figure caption (auxiliary; use the image as primary evidence)

The image depicts a diagram with two charged particles and a point in space. 

- On the top left, there is a red circle with a plus sign labeled "3.0 nC," indicating a positive charge of 3.0 nanocoulombs.
- Directly below it, separated by 10 cm, is a teal circle with a minus sign labeled "-3.0 nC," indicating a negative charge of -3.0 nanocoulombs.
- A point in space marked by a small black dot is located 5.0 cm to the right of the positive charge.
- The distances "5.0 cm" (horizontal) and "10 cm" (vertical) are indicated with double-headed arrows.
- The overall layout shows a right angle between the distance from the positive charge to the point and the vertical distance between the charges.

This diagram appears to represent an electrostatic configuration for analyzing electric fields or forces.

\#\# Dataset metadata

Electrostatics; Physical Model Grounding Reasoning, Spatial Relation Reasoning

\paragraph{Recorded answer/context.}
C: C: 1.0 \textbackslash{}times 10\textasciicircum{}\{4\}N/C

\paragraph{Figure/image path.}
/root/proposal\_for\_physic/science-mango/phyx\_mini\_run/phyx\_data/test\_image/915.png

\paragraph{Formalization target.}
create a compiling Lean file with sorry bodies at `PhyXMiniProblems/problem\_phyx\_mini\_0915.lean`. The Lean declarations must preserve the physical quantities, dimensions or dimensional roles, figure labels, governing-law hypotheses, and final relation expressed by this problem.
Use Mathlib/Physlib names found through LeanExplore where available. If a physics API is missing, introduce faithful local abstractions rather than scalar placeholder aliases.

\begin{theorem}[Physics formalization target]
\label{thm:physics:phyx_mini_0915:target}
\lean{PhyXMiniProblems.ProblemPhyXMini0915.problem_phyx_mini_0915}
\uses{def:physics:phyx-mini-0915:phyxminiproblems-problemphyxmini0915-dipolefieldsetup, def:physics:phyx-mini-0915:phyxminiproblems-problemphyxmini0915-matcheswrittenelectrostaticscenario, def:physics:phyx-mini-0915:phyxminiproblems-problemphyxmini0915-matchessupplieddipolefieldfigure, def:physics:phyx-mini-0915:phyxminiproblems-problemphyxmini0915-usesdisplayedrightanglegeometry, def:physics:phyx-mini-0915:phyxminiproblems-problemphyxmini0915-hasphysicalelectrostaticparameters, def:physics:phyx-mini-0915:phyxminiproblems-problemphyxmini0915-usestextbookvacuumcoulombconstant, def:physics:phyx-mini-0915:phyxminiproblems-problemphyxmini0915-satisfiespointchargecoulomblawandsuperposition, def:physics:phyx-mini-0915:phyxminiproblems-problemphyxmini0915-resultantfieldstrengthinnewtonspercoulomb, def:physics:phyx-mini-0915:phyxminiproblems-problemphyxmini0915-recordeddatasetanswer, def:physics:phyx-mini-0915:phyxminiproblems-problemphyxmini0915-roundstonearestresolution, def:physics:phyx-mini-0915:phyxminiproblems-problemphyxmini0915-isuniquematchinganswer}
The assigned autoformalize agent should translate this physics problem into Lean declarations in the covered file, with theorem and lemma proof bodies written as `by sorry`.
\end{theorem}
\begin{proof}
This is an autoformalization task, not a proof task. Produce statements that can later be proved by the physics prover without weakening the physical model.
\end{proof}

% --- Archon declaration topology begin ---
\section{Declaration topology}
\begin{definition}[coulomb Constant Dimension]
  \label{def:physics:phyx-mini-0915:phyxminiproblems-problemphyxmini0915-coulombconstantdimension}
  \lean{PhyXMiniProblems.ProblemPhyXMini0915.coulombConstantDimension}
  Coulomb's constant has physical dimension M L\textasciicircum{}3 T\textasciicircum{}-2 C\textasciicircum{}-2
\end{definition}
\begin{proof}
  This is the declared model definition of coulomb Constant Dimension.
\end{proof}
\begin{definition}[electric Field Dimension]
  \label{def:physics:phyx-mini-0915:phyxminiproblems-problemphyxmini0915-electricfielddimension}
  \lean{PhyXMiniProblems.ProblemPhyXMini0915.electricFieldDimension}
  Electric-field strength has physical dimension M L T\textasciicircum{}-2 C\textasciicircum{}-1
\end{definition}
\begin{proof}
  This is the declared model definition of electric Field Dimension.
\end{proof}
\begin{definition}[Length Quantity]
  \label{def:physics:phyx-mini-0915:phyxminiproblems-problemphyxmini0915-lengthquantity}
  \lean{PhyXMiniProblems.ProblemPhyXMini0915.LengthQuantity}
  A nonnegative physical length, independent of its readout unit
\end{definition}
\begin{proof}
  This is the declared model definition of Length Quantity.
\end{proof}
\begin{definition}[Signed Charge Quantity]
  \label{def:physics:phyx-mini-0915:phyxminiproblems-problemphyxmini0915-signedchargequantity}
  \lean{PhyXMiniProblems.ProblemPhyXMini0915.SignedChargeQuantity}
  A signed physical electric charge
\end{definition}
\begin{proof}
  This is the declared model definition of Signed Charge Quantity.
\end{proof}
\begin{definition}[Planar Position Quantity]
  \label{def:physics:phyx-mini-0915:phyxminiproblems-problemphyxmini0915-planarpositionquantity}
  \lean{PhyXMiniProblems.ProblemPhyXMini0915.PlanarPositionQuantity}
  A physical position vector in the plane of the figure
\end{definition}
\begin{proof}
  This is the declared model definition of Planar Position Quantity.
\end{proof}
\begin{definition}[Planar Electric Field Quantity]
  \label{def:physics:phyx-mini-0915:phyxminiproblems-problemphyxmini0915-planarelectricfieldquantity}
  \lean{PhyXMiniProblems.ProblemPhyXMini0915.PlanarElectricFieldQuantity}
  \uses{def:physics:phyx-mini-0915:phyxminiproblems-problemphyxmini0915-electricfielddimension}
  A physical electric-field vector in the plane of the figure
\end{definition}
\begin{proof}
  This is the declared model definition of Planar Electric Field Quantity.
\end{proof}
\begin{definition}[Coulomb Constant Quantity]
  \label{def:physics:phyx-mini-0915:phyxminiproblems-problemphyxmini0915-coulombconstantquantity}
  \lean{PhyXMiniProblems.ProblemPhyXMini0915.CoulombConstantQuantity}
  \uses{def:physics:phyx-mini-0915:phyxminiproblems-problemphyxmini0915-coulombconstantdimension}
  A nonnegative dimensionful value of Coulomb's constant
\end{definition}
\begin{proof}
  This is the declared model definition of Coulomb Constant Quantity.
\end{proof}
\begin{definition}[length Readout]
  \label{def:physics:phyx-mini-0915:phyxminiproblems-problemphyxmini0915-lengthreadout}
  \lean{PhyXMiniProblems.ProblemPhyXMini0915.lengthReadout}
  \uses{def:physics:phyx-mini-0915:phyxminiproblems-problemphyxmini0915-lengthquantity}
  Read a physical length in a selected Physlib length unit
\end{definition}
\begin{proof}
  This is the declared model definition of length Readout.
\end{proof}
\begin{definition}[length In Meters]
  \label{def:physics:phyx-mini-0915:phyxminiproblems-problemphyxmini0915-lengthinmeters}
  \lean{PhyXMiniProblems.ProblemPhyXMini0915.lengthInMeters}
  \uses{def:physics:phyx-mini-0915:phyxminiproblems-problemphyxmini0915-lengthquantity, def:physics:phyx-mini-0915:phyxminiproblems-problemphyxmini0915-lengthreadout}
  Read a physical length in coherent-SI metres
\end{definition}
\begin{proof}
  This is the declared model definition of length In Meters.
\end{proof}
\begin{definition}[length In Centimeters]
  \label{def:physics:phyx-mini-0915:phyxminiproblems-problemphyxmini0915-lengthincentimeters}
  \lean{PhyXMiniProblems.ProblemPhyXMini0915.lengthInCentimeters}
  \uses{def:physics:phyx-mini-0915:phyxminiproblems-problemphyxmini0915-lengthquantity, def:physics:phyx-mini-0915:phyxminiproblems-problemphyxmini0915-lengthreadout}
  Read a physical length in centimetres
\end{definition}
\begin{proof}
  This is the declared model definition of length In Centimeters.
\end{proof}
\begin{definition}[charge In Coulombs]
  \label{def:physics:phyx-mini-0915:phyxminiproblems-problemphyxmini0915-chargeincoulombs}
  \lean{PhyXMiniProblems.ProblemPhyXMini0915.chargeInCoulombs}
  \uses{def:physics:phyx-mini-0915:phyxminiproblems-problemphyxmini0915-signedchargequantity}
  Read a signed physical charge in coherent-SI coulombs
\end{definition}
\begin{proof}
  This is the declared model definition of charge In Coulombs.
\end{proof}
\begin{definition}[charge In Nanocoulombs]
  \label{def:physics:phyx-mini-0915:phyxminiproblems-problemphyxmini0915-chargeinnanocoulombs}
  \lean{PhyXMiniProblems.ProblemPhyXMini0915.chargeInNanocoulombs}
  \uses{def:physics:phyx-mini-0915:phyxminiproblems-problemphyxmini0915-signedchargequantity, def:physics:phyx-mini-0915:phyxminiproblems-problemphyxmini0915-chargeincoulombs}
  Read a signed physical charge in nanocoulombs
\end{definition}
\begin{proof}
  This is the declared model definition of charge In Nanocoulombs.
\end{proof}
\begin{definition}[planar Position In Meters]
  \label{def:physics:phyx-mini-0915:phyxminiproblems-problemphyxmini0915-planarpositioninmeters}
  \lean{PhyXMiniProblems.ProblemPhyXMini0915.planarPositionInMeters}
  \uses{def:physics:phyx-mini-0915:phyxminiproblems-problemphyxmini0915-planarpositionquantity}
  Read a planar physical position in coherent-SI metres
\end{definition}
\begin{proof}
  This is the declared model definition of planar Position In Meters.
\end{proof}
\begin{definition}[position As Physlib Space]
  \label{def:physics:phyx-mini-0915:phyxminiproblems-problemphyxmini0915-positionasphyslibspace}
  \lean{PhyXMiniProblems.ProblemPhyXMini0915.positionAsPhyslibSpace}
  \uses{def:physics:phyx-mini-0915:phyxminiproblems-problemphyxmini0915-planarpositionquantity, def:physics:phyx-mini-0915:phyxminiproblems-problemphyxmini0915-planarpositioninmeters}
  Regard a dimensionful planar position's SI readout as a Physlib point
\end{definition}
\begin{proof}
  This is the declared model definition of position As Physlib Space.
\end{proof}
\begin{definition}[planar Electric Field In Newtons Per Coulomb]
  \label{def:physics:phyx-mini-0915:phyxminiproblems-problemphyxmini0915-planarelectricfieldinnewtonspercoulomb}
  \lean{PhyXMiniProblems.ProblemPhyXMini0915.planarElectricFieldInNewtonsPerCoulomb}
  \uses{def:physics:phyx-mini-0915:phyxminiproblems-problemphyxmini0915-planarelectricfieldquantity}
  Read a physical planar electric-field vector in newtons per coulomb
\end{definition}
\begin{proof}
  This is the declared model definition of planar Electric Field In Newtons Per Coulomb.
\end{proof}
\begin{definition}[coulomb Constant In Newton Meters Squared Per Coulomb Squared]
  \label{def:physics:phyx-mini-0915:phyxminiproblems-problemphyxmini0915-coulombconstantinnewtonmeterssquaredpercoulombsquared}
  \lean{PhyXMiniProblems.ProblemPhyXMini0915.coulombConstantInNewtonMetersSquaredPerCoulombSquared}
  \uses{def:physics:phyx-mini-0915:phyxminiproblems-problemphyxmini0915-coulombconstantquantity}
  Read Coulomb's constant in coherent SI N m\textasciicircum{}2 / C\textasciicircum{}2
\end{definition}
\begin{proof}
  This is the declared model definition of coulomb Constant In Newton Meters Squared Per Coulomb Squared.
\end{proof}
\begin{definition}[Charge Source]
  \label{def:physics:phyx-mini-0915:phyxminiproblems-problemphyxmini0915-chargesource}
  \lean{PhyXMiniProblems.ProblemPhyXMini0915.ChargeSource}
  The two source charges distinguished by their vertical positions
\end{definition}
\begin{proof}
  This is the declared model definition of Charge Source.
\end{proof}
\begin{definition}[Figure Location]
  \label{def:physics:phyx-mini-0915:phyxminiproblems-problemphyxmini0915-figurelocation}
  \lean{PhyXMiniProblems.ProblemPhyXMini0915.FigureLocation}
  All three marked positions in the supplied raster
\end{definition}
\begin{proof}
  This is the declared model definition of Figure Location.
\end{proof}
\begin{definition}[figure Location]
  \label{def:physics:phyx-mini-0915:phyxminiproblems-problemphyxmini0915-chargesource-figurelocation}
  \lean{PhyXMiniProblems.ProblemPhyXMini0915.ChargeSource.figureLocation}
  \uses{def:physics:phyx-mini-0915:phyxminiproblems-problemphyxmini0915-chargesource, def:physics:phyx-mini-0915:phyxminiproblems-problemphyxmini0915-figurelocation}
  Regard a source label as the corresponding marked figure location
\end{definition}
\begin{proof}
  This is the declared model definition of figure Location.
\end{proof}
\begin{definition}[Distance Arrow]
  \label{def:physics:phyx-mini-0915:phyxminiproblems-problemphyxmini0915-distancearrow}
  \lean{PhyXMiniProblems.ProblemPhyXMini0915.DistanceArrow}
  The two double-headed distance arrows shown in the raster
\end{definition}
\begin{proof}
  This is the declared model definition of Distance Arrow.
\end{proof}
\begin{definition}[Marker Kind]
  \label{def:physics:phyx-mini-0915:phyxminiproblems-problemphyxmini0915-markerkind}
  \lean{PhyXMiniProblems.ProblemPhyXMini0915.MarkerKind}
  The visual role of a marker at a displayed location
\end{definition}
\begin{proof}
  This is the declared model definition of Marker Kind.
\end{proof}
\begin{definition}[Marker Color]
  \label{def:physics:phyx-mini-0915:phyxminiproblems-problemphyxmini0915-markercolor}
  \lean{PhyXMiniProblems.ProblemPhyXMini0915.MarkerColor}
  Marker colors distinguished in the primary image
\end{definition}
\begin{proof}
  This is the declared model definition of Marker Color.
\end{proof}
\begin{definition}[Charge Sign Glyph]
  \label{def:physics:phyx-mini-0915:phyxminiproblems-problemphyxmini0915-chargesignglyph}
  \lean{PhyXMiniProblems.ProblemPhyXMini0915.ChargeSignGlyph}
  The sign glyph printed within a charged circle
\end{definition}
\begin{proof}
  This is the declared model definition of Charge Sign Glyph.
\end{proof}
\begin{definition}[Arrow Orientation]
  \label{def:physics:phyx-mini-0915:phyxminiproblems-problemphyxmini0915-arroworientation}
  \lean{PhyXMiniProblems.ProblemPhyXMini0915.ArrowOrientation}
  The qualitative direction of a displayed distance arrow
\end{definition}
\begin{proof}
  This is the declared model definition of Arrow Orientation.
\end{proof}
\begin{definition}[Electrostatic Source Model]
  \label{def:physics:phyx-mini-0915:phyxminiproblems-problemphyxmini0915-electrostaticsourcemodel}
  \lean{PhyXMiniProblems.ProblemPhyXMini0915.ElectrostaticSourceModel}
  The source idealization used in the elementary electrostatic model
\end{definition}
\begin{proof}
  This is the declared model definition of Electrostatic Source Model.
\end{proof}
\begin{definition}[Electrostatic Medium]
  \label{def:physics:phyx-mini-0915:phyxminiproblems-problemphyxmini0915-electrostaticmedium}
  \lean{PhyXMiniProblems.ProblemPhyXMini0915.ElectrostaticMedium}
  The medium surrounding the source charges and observation dot
\end{definition}
\begin{proof}
  This is the declared model definition of Electrostatic Medium.
\end{proof}
\begin{definition}[Dipole Field Figure]
  \label{def:physics:phyx-mini-0915:phyxminiproblems-problemphyxmini0915-dipolefieldfigure}
  \lean{PhyXMiniProblems.ProblemPhyXMini0915.DipoleFieldFigure}
  \uses{def:physics:phyx-mini-0915:phyxminiproblems-problemphyxmini0915-chargesource, def:physics:phyx-mini-0915:phyxminiproblems-problemphyxmini0915-figurelocation, def:physics:phyx-mini-0915:phyxminiproblems-problemphyxmini0915-distancearrow, def:physics:phyx-mini-0915:phyxminiproblems-problemphyxmini0915-markerkind, def:physics:phyx-mini-0915:phyxminiproblems-problemphyxmini0915-markercolor, def:physics:phyx-mini-0915:phyxminiproblems-problemphyxmini0915-chargesignglyph, def:physics:phyx-mini-0915:phyxminiproblems-problemphyxmini0915-arroworientation}
  Literal presentation information from image 915.png. Its real fields are unit-labelled printed readouts, not replacements for physical quantities. There is deliberately no field-strength value in this figure structure
\end{definition}
\begin{proof}
  This is the declared model definition of Dipole Field Figure.
\end{proof}
\begin{definition}[Point Charge Source]
  \label{def:physics:phyx-mini-0915:phyxminiproblems-problemphyxmini0915-pointchargesource}
  \lean{PhyXMiniProblems.ProblemPhyXMini0915.PointChargeSource}
  \uses{def:physics:phyx-mini-0915:phyxminiproblems-problemphyxmini0915-signedchargequantity, def:physics:phyx-mini-0915:phyxminiproblems-problemphyxmini0915-planarpositionquantity}
  A charged source retains both its dimensionful charge and dimensionful planar position. This is a physical point-charge object, rather than a transparent scalar alias for charge
\end{definition}
\begin{proof}
  This is the declared model definition of Point Charge Source.
\end{proof}
\begin{definition}[Dipole Field Setup]
  \label{def:physics:phyx-mini-0915:phyxminiproblems-problemphyxmini0915-dipolefieldsetup}
  \lean{PhyXMiniProblems.ProblemPhyXMini0915.DipoleFieldSetup}
  \uses{def:physics:phyx-mini-0915:phyxminiproblems-problemphyxmini0915-lengthquantity, def:physics:phyx-mini-0915:phyxminiproblems-problemphyxmini0915-planarpositionquantity, def:physics:phyx-mini-0915:phyxminiproblems-problemphyxmini0915-planarelectricfieldquantity, def:physics:phyx-mini-0915:phyxminiproblems-problemphyxmini0915-coulombconstantquantity, def:physics:phyx-mini-0915:phyxminiproblems-problemphyxmini0915-chargesource, def:physics:phyx-mini-0915:phyxminiproblems-problemphyxmini0915-electrostaticsourcemodel, def:physics:phyx-mini-0915:phyxminiproblems-problemphyxmini0915-electrostaticmedium, def:physics:phyx-mini-0915:phyxminiproblems-problemphyxmini0915-dipolefieldfigure, def:physics:phyx-mini-0915:phyxminiproblems-problemphyxmini0915-pointchargesource}
  Independent physical data for the two-charge configuration. The individual field contributions and their resultant are independent quantities. Neither is manufactured from an answer value
\end{definition}
\begin{proof}
  This is the declared model definition of Dipole Field Setup.
\end{proof}
\begin{definition}[position At]
  \label{def:physics:phyx-mini-0915:phyxminiproblems-problemphyxmini0915-dipolefieldsetup-positionat}
  \lean{PhyXMiniProblems.ProblemPhyXMini0915.DipoleFieldSetup.positionAt}
  \uses{def:physics:phyx-mini-0915:phyxminiproblems-problemphyxmini0915-planarpositionquantity, def:physics:phyx-mini-0915:phyxminiproblems-problemphyxmini0915-figurelocation, def:physics:phyx-mini-0915:phyxminiproblems-problemphyxmini0915-dipolefieldsetup}
  Physical position of each marked location
\end{definition}
\begin{proof}
  This is the declared model definition of position At.
\end{proof}
\begin{definition}[Matches Written Electrostatic Scenario]
  \label{def:physics:phyx-mini-0915:phyxminiproblems-problemphyxmini0915-matcheswrittenelectrostaticscenario}
  \lean{PhyXMiniProblems.ProblemPhyXMini0915.MatchesWrittenElectrostaticScenario}
  \uses{def:physics:phyx-mini-0915:phyxminiproblems-problemphyxmini0915-dipolefieldsetup}
  The two sources are stationary point charges in vacuum
\end{definition}
\begin{proof}
  This is the declared model definition of Matches Written Electrostatic Scenario.
\end{proof}
\begin{definition}[Matches Supplied Dipole Field Figure]
  \label{def:physics:phyx-mini-0915:phyxminiproblems-problemphyxmini0915-matchessupplieddipolefieldfigure}
  \lean{PhyXMiniProblems.ProblemPhyXMini0915.MatchesSuppliedDipoleFieldFigure}
  \uses{def:physics:phyx-mini-0915:phyxminiproblems-problemphyxmini0915-lengthincentimeters, def:physics:phyx-mini-0915:phyxminiproblems-problemphyxmini0915-chargeinnanocoulombs, def:physics:phyx-mini-0915:phyxminiproblems-problemphyxmini0915-dipolefieldsetup}
  Primary-raster evidence and calibration. The upper red charge is +3.0 nC, the lower teal charge is -3.0 nC, the dot is black, and the horizontal and vertical arrows read 5.0 cm and 10 cm. No resultant-field value or answer label occurs in this structure
\end{definition}
\begin{proof}
  This is the declared model definition of Matches Supplied Dipole Field Figure.
\end{proof}
\begin{definition}[Uses Displayed Right Angle Geometry]
  \label{def:physics:phyx-mini-0915:phyxminiproblems-problemphyxmini0915-usesdisplayedrightanglegeometry}
  \lean{PhyXMiniProblems.ProblemPhyXMini0915.UsesDisplayedRightAngleGeometry}
  \uses{def:physics:phyx-mini-0915:phyxminiproblems-problemphyxmini0915-lengthinmeters, def:physics:phyx-mini-0915:phyxminiproblems-problemphyxmini0915-planarpositioninmeters, def:physics:phyx-mini-0915:phyxminiproblems-problemphyxmini0915-dipolefieldsetup}
  Cartesian realization of the displayed geometry. The origin is chosen at the upper positive charge, the positive x-axis points toward the dot, and the positive y-axis points upward. Thus the lower negative charge has vertical coordinate -10 cm. These equations express only figure geometry
\end{definition}
\begin{proof}
  This is the declared model definition of Uses Displayed Right Angle Geometry.
\end{proof}
\begin{definition}[Has Physical Electrostatic Parameters]
  \label{def:physics:phyx-mini-0915:phyxminiproblems-problemphyxmini0915-hasphysicalelectrostaticparameters}
  \lean{PhyXMiniProblems.ProblemPhyXMini0915.HasPhysicalElectrostaticParameters}
  \uses{def:physics:phyx-mini-0915:phyxminiproblems-problemphyxmini0915-lengthinmeters, def:physics:phyx-mini-0915:phyxminiproblems-problemphyxmini0915-chargeincoulombs, def:physics:phyx-mini-0915:phyxminiproblems-problemphyxmini0915-planarpositioninmeters, def:physics:phyx-mini-0915:phyxminiproblems-problemphyxmini0915-coulombconstantinnewtonmeterssquaredpercoulombsquared, def:physics:phyx-mini-0915:phyxminiproblems-problemphyxmini0915-dipolefieldsetup}
  Positivity and non-collision conditions for the physical configuration
\end{definition}
\begin{proof}
  This is the declared model definition of Has Physical Electrostatic Parameters.
\end{proof}
\begin{definition}[Uses Textbook Vacuum Coulomb Constant]
  \label{def:physics:phyx-mini-0915:phyxminiproblems-problemphyxmini0915-usestextbookvacuumcoulombconstant}
  \lean{PhyXMiniProblems.ProblemPhyXMini0915.UsesTextbookVacuumCoulombConstant}
  \uses{def:physics:phyx-mini-0915:phyxminiproblems-problemphyxmini0915-coulombconstantinnewtonmeterssquaredpercoulombsquared, def:physics:phyx-mini-0915:phyxminiproblems-problemphyxmini0915-dipolefieldsetup}
  The dimensionful Coulomb constant is linked to Physlib's electromagnetic system and calibrated to the rounded textbook value 9.0 * 10\textasciicircum{}9 N m\textasciicircum{}2/C\textasciicircum{}2. This contains no field-strength answer
\end{definition}
\begin{proof}
  This is the declared model definition of Uses Textbook Vacuum Coulomb Constant.
\end{proof}
\begin{definition}[Satisfies Point Charge Coulomb Law And Superposition]
  \label{def:physics:phyx-mini-0915:phyxminiproblems-problemphyxmini0915-satisfiespointchargecoulomblawandsuperposition}
  \lean{PhyXMiniProblems.ProblemPhyXMini0915.SatisfiesPointChargeCoulombLawAndSuperposition}
  \uses{def:physics:phyx-mini-0915:phyxminiproblems-problemphyxmini0915-chargeincoulombs, def:physics:phyx-mini-0915:phyxminiproblems-problemphyxmini0915-planarpositioninmeters, def:physics:phyx-mini-0915:phyxminiproblems-problemphyxmini0915-positionasphyslibspace, def:physics:phyx-mini-0915:phyxminiproblems-problemphyxmini0915-planarelectricfieldinnewtonspercoulomb, def:physics:phyx-mini-0915:phyxminiproblems-problemphyxmini0915-coulombconstantinnewtonmeterssquaredpercoulombsquared, def:physics:phyx-mini-0915:phyxminiproblems-problemphyxmini0915-chargesource, def:physics:phyx-mini-0915:phyxminiproblems-problemphyxmini0915-dipolefieldsetup}
  The vector point-charge law is E = k q (r - r\_source) / ‖r - r\_source‖\textasciicircum{}3. The second field states linear superposition, and the third connects the independent dimensionful resultant to Physlib's electric field at the observation event. These are general laws and contain no numerical resultant, rounding interval, or answer choice
\end{definition}
\begin{proof}
  This is the declared model definition of Satisfies Point Charge Coulomb Law And Superposition.
\end{proof}
\begin{definition}[resultant Field Strength In Newtons Per Coulomb]
  \label{def:physics:phyx-mini-0915:phyxminiproblems-problemphyxmini0915-resultantfieldstrengthinnewtonspercoulomb}
  \lean{PhyXMiniProblems.ProblemPhyXMini0915.resultantFieldStrengthInNewtonsPerCoulomb}
  \uses{def:physics:phyx-mini-0915:phyxminiproblems-problemphyxmini0915-planarelectricfieldinnewtonspercoulomb, def:physics:phyx-mini-0915:phyxminiproblems-problemphyxmini0915-dipolefieldsetup}
  Magnitude of the resultant field at the dot, in newtons per coulomb
\end{definition}
\begin{proof}
  This is the declared model definition of resultant Field Strength In Newtons Per Coulomb.
\end{proof}
\begin{definition}[Answer Choice]
  \label{def:physics:phyx-mini-0915:phyxminiproblems-problemphyxmini0915-answerchoice}
  \lean{PhyXMiniProblems.ProblemPhyXMini0915.AnswerChoice}
  Labels of the four answer choices supplied with the question
\end{definition}
\begin{proof}
  This is the declared model definition of Answer Choice.
\end{proof}
\begin{definition}[displayed Field Strength In Newtons Per Coulomb]
  \label{def:physics:phyx-mini-0915:phyxminiproblems-problemphyxmini0915-answerchoice-displayedfieldstrengthinnewtonspercoulomb}
  \lean{PhyXMiniProblems.ProblemPhyXMini0915.AnswerChoice.displayedFieldStrengthInNewtonsPerCoulomb}
  \uses{def:physics:phyx-mini-0915:phyxminiproblems-problemphyxmini0915-answerchoice}
  The field strength printed for each choice, in newtons per coulomb
\end{definition}
\begin{proof}
  This is the declared model definition of displayed Field Strength In Newtons Per Coulomb.
\end{proof}
\begin{definition}[displayed Resolution In Newtons Per Coulomb]
  \label{def:physics:phyx-mini-0915:phyxminiproblems-problemphyxmini0915-answerchoice-displayedresolutioninnewtonspercoulomb}
  \lean{PhyXMiniProblems.ProblemPhyXMini0915.AnswerChoice.displayedResolutionInNewtonsPerCoulomb}
  \uses{def:physics:phyx-mini-0915:phyxminiproblems-problemphyxmini0915-answerchoice}
  Resolution implied by the two significant digits printed in each choice
\end{definition}
\begin{proof}
  This is the declared model definition of displayed Resolution In Newtons Per Coulomb.
\end{proof}
\begin{definition}[recorded Dataset Answer]
  \label{def:physics:phyx-mini-0915:phyxminiproblems-problemphyxmini0915-recordeddatasetanswer}
  \lean{PhyXMiniProblems.ProblemPhyXMini0915.recordedDatasetAnswer}
  \uses{def:physics:phyx-mini-0915:phyxminiproblems-problemphyxmini0915-answerchoice}
  The answer label recorded by the source dataset
\end{definition}
\begin{proof}
  This is the declared model definition of recorded Dataset Answer.
\end{proof}
\begin{definition}[Rounds To Nearest Resolution]
  \label{def:physics:phyx-mini-0915:phyxminiproblems-problemphyxmini0915-roundstonearestresolution}
  \lean{PhyXMiniProblems.ProblemPhyXMini0915.RoundsToNearestResolution}
  A value rounds to a display at the stated positive resolution
\end{definition}
\begin{proof}
  This is the declared model definition of Rounds To Nearest Resolution.
\end{proof}
\begin{definition}[Answer Matches Resultant Field Strength]
  \label{def:physics:phyx-mini-0915:phyxminiproblems-problemphyxmini0915-answermatchesresultantfieldstrength}
  \lean{PhyXMiniProblems.ProblemPhyXMini0915.AnswerMatchesResultantFieldStrength}
  \uses{def:physics:phyx-mini-0915:phyxminiproblems-problemphyxmini0915-dipolefieldsetup, def:physics:phyx-mini-0915:phyxminiproblems-problemphyxmini0915-resultantfieldstrengthinnewtonspercoulomb, def:physics:phyx-mini-0915:phyxminiproblems-problemphyxmini0915-answerchoice, def:physics:phyx-mini-0915:phyxminiproblems-problemphyxmini0915-roundstonearestresolution}
  A choice displays the resultant field strength at its printed precision
\end{definition}
\begin{proof}
  This is the declared model definition of Answer Matches Resultant Field Strength.
\end{proof}
\begin{definition}[Is Unique Matching Answer]
  \label{def:physics:phyx-mini-0915:phyxminiproblems-problemphyxmini0915-isuniquematchinganswer}
  \lean{PhyXMiniProblems.ProblemPhyXMini0915.IsUniqueMatchingAnswer}
  \uses{def:physics:phyx-mini-0915:phyxminiproblems-problemphyxmini0915-dipolefieldsetup, def:physics:phyx-mini-0915:phyxminiproblems-problemphyxmini0915-answerchoice, def:physics:phyx-mini-0915:phyxminiproblems-problemphyxmini0915-answermatchesresultantfieldstrength}
  A choice is the unique displayed answer matching the modeled strength
\end{definition}
\begin{proof}
  This is the declared model definition of Is Unique Matching Answer.
\end{proof}
\begin{lemma}[resultant Field Strength rounding bounds]
  \label{lem:physics:phyx-mini-0915:phyxminiproblems-problemphyxmini0915-resultantfieldstrength-rounding-bounds}
  \lean{PhyXMiniProblems.ProblemPhyXMini0915.resultantFieldStrength_rounding_bounds}
  \uses{def:physics:phyx-mini-0915:phyxminiproblems-problemphyxmini0915-dipolefieldsetup, def:physics:phyx-mini-0915:phyxminiproblems-problemphyxmini0915-matchessupplieddipolefieldfigure, def:physics:phyx-mini-0915:phyxminiproblems-problemphyxmini0915-usesdisplayedrightanglegeometry, def:physics:phyx-mini-0915:phyxminiproblems-problemphyxmini0915-hasphysicalelectrostaticparameters, def:physics:phyx-mini-0915:phyxminiproblems-problemphyxmini0915-usestextbookvacuumcoulombconstant, def:physics:phyx-mini-0915:phyxminiproblems-problemphyxmini0915-satisfiespointchargecoulomblawandsuperposition, def:physics:phyx-mini-0915:phyxminiproblems-problemphyxmini0915-resultantfieldstrengthinnewtonspercoulomb}
  The two vector contributions give a resultant strength between 9500 and 10500 N/C. This is a derived numerical statement, not a premise
\end{lemma}
\begin{proof}
  The conclusion follows from the governing relations and figure readouts named in the statement of resultant Field Strength rounding bounds.
\end{proof}
% --- Archon declaration topology end ---
% --- Archon physics formalization source end ---
